\documentclass[11pt]{amsart}
\usepackage[margin=1.1in]{geometry}
\usepackage{amsmath,amssymb,amsthm}
\usepackage{booktabs}
\usepackage{array}
\usepackage{microtype}
\usepackage[colorlinks=true,linkcolor=blue,citecolor=blue,urlcolor=blue]{hyperref}

\newtheorem{theorem}{Theorem}[section]
\newtheorem{lemma}[theorem]{Lemma}
\newtheorem{proposition}[theorem]{Proposition}
\newtheorem{corollary}[theorem]{Corollary}
\newtheorem{introthm}{Theorem}
\renewcommand{\theintrothm}{\Alph{introthm}}
\newtheorem{introcor}[introthm]{Corollary}
\theoremstyle{definition}
\newtheorem{definition}[theorem]{Definition}
\newtheorem{question}[theorem]{Question}
\theoremstyle{remark}
\newtheorem{remark}[theorem]{Remark}

\newcommand{\Z}{\mathbb{Z}}
\newcommand{\R}{\mathbb{R}}
\newcommand{\C}{\mathbb{C}}
\newcommand{\T}{\mathbb{T}}
\newcommand{\D}{\mathbb{D}}
\newcommand{\HH}{\mathbb{H}}
\newcommand{\ip}[2]{\langle #1,#2\rangle}
\newcommand{\tr}{\operatorname{tr}}
\newcommand{\Tr}{\operatorname{Tr}}
\renewcommand{\Re}{\operatorname{Re}}
\renewcommand{\Im}{\operatorname{Im}}
\newcommand{\FD}{F_{\mathrm{DKZ}}}
\newcommand{\IME}{I_{\mathrm{ME}}}
\newcommand{\Qc}{Q_{\mathrm{c}}}
\newcommand{\kb}{\bar\kappa}
\newcommand{\sprod}{\mathrm{sc}}
\newcommand{\fw}{\mathrm{fw}}
\newcommand{\one}{\mathbf{1}}
\setlength{\emergencystretch}{3em}

\begin{document}

\title[Discrete Hilbert transform rearrangement and CGLMP]{A rearrangement inequality for the discrete Hilbert
transform, with an application to optimal CGLMP measurements}

\author{Ansh Mishra}
\author{Aryan Senthilkumar}

\date{September 2026}

\begin{abstract}
For disjoint subsets $A,B$ of the cyclic group $\Z_N$ of prescribed sizes we show that the pairing
$\sum_{x\in A,\,y\in B}\cot(\pi(x-y)/N)$, the bilinear form of the discrete Hilbert transform, is largest
exactly when $A$ and $B$ are adjacent intervals with $B$ immediately before $A$. The proof has three
ingredients: an exact exchange identity, which forces every maximiser to be cyclically ordered; a continuous
version of the inequality, based on the circle analogue of a theorem of Stein and Weiss (the distribution of the
conjugate function of an indicator function depends only on the measure of the set); and an explicit estimate of
the discretisation error at the junctions between blocks, which excludes cyclically ordered configurations that
wind around the circle more than once.

We apply the inequality to Open Quantum Problem~27B of IQOQI Vienna, which asks to show that the measurements
of Durt, Kaszlikowski and \.Zukowski (DKZ) are necessarily the optimal ones for the CGLMP$_d$ Bell inequality on
maximally entangled states. The optimality statement is equivalent to an inequality for $d$-tuples of unitaries of
order four (a $\Z_4$ clock model). We prove that inequality for every $d$ when the unitaries commute. In physical
terms: among maximally entangled strategies whose four link operators coincide, the DKZ measurements are optimal
and, up to local unitaries, unique, in every dimension $d$. We also prove the inequality, for every $d$, for a class
of genuinely non-commuting families in which each unitary switches between the values of two optimal
configurations. The general non-commuting case remains open for $d\ge 21$; for $d\le 20$ it is settled by
exact computer-assisted certificates.
\end{abstract}

\maketitle

\section{Introduction}\label{sec:intro}

\subsection{The problem}
The CGLMP inequalities of Collins, Gisin, Linden, Massar and Popescu~\cite{CGLMP} are Bell inequalities for two
parties with two measurement settings each and $d$ outcomes per measurement. Write $A_1,A_2$ and $B_1,B_2$ for
the outcomes (elements of $\Z_d$) of Alice's and Bob's measurements. The CGLMP expression is
\begin{multline}\label{eq:cglmp}
I_d=\sum_{k=0}^{\lfloor d/2\rfloor-1}\Bigl(1-\frac{2k}{d-1}\Bigr)\Bigl(\bigl[P(A_1=B_1+k)+P(B_1=A_2+k+1)\\
+P(A_2=B_2+k)+P(B_2=A_1+k)\bigr]
-\bigl[P(A_1=B_1-k-1)+P(B_1=A_2-k)\\
+P(A_2=B_2-k-1)+P(B_2=A_1-k-1)\bigr]\Bigr),
\end{multline}
where $P(A_x=B_y+k)$ is the probability that $A_x-B_y\equiv k\pmod d$ when the settings $x$ and $y$ are
used.\footnote{This is the reading under which the local bound is derived in \cite[eqs.~(7)--(11)]{CGLMP} and
under which the measurements of \cite[eqs.~(12)--(15)]{CGLMP} give the value \eqref{eq:IME}; see check (R7) in
Section~\ref{sec:numerics}.} Local hidden variable models satisfy $I_d\le 2$~\cite{CGLMP}. On the maximally
entangled state of two $d$-level systems, the measurements found numerically by Kaszlikowski et
al.~\cite{KGZMZ} and described analytically by Durt, Kaszlikowski and \.Zukowski~\cite{DKZ}, namely diagonal
phase shifts followed by a discrete Fourier transform, which we call the \emph{DKZ measurements}, give the value
\begin{equation}\label{eq:IME}
\IME(d)=\frac{4}{d(d-1)}\sum_{j=1}^{d-1}(d-j)\sec\frac{\pi j}{2d},
\end{equation}
for example $\IME(3)=2.8729\ldots$ and $\IME(4)=2.8962\ldots$. Numerical searches~\cite{KGZMZ,DKZ,CGLMP} found
no better measurements on the maximally entangled state, while for $d\ge3$ the maximally entangled state is not
the optimal state for $I_d$~\cite{ADGL}. Part~B of Open Quantum Problem~27 of the IQOQI Vienna list (``The power
of CGLMP inequalities'', \cite{OQP}) records that the observables which numerically maximise the CGLMP violation on
a maximally entangled state are measurements in the computational basis transformed by a discrete Fourier
transform and diagonal unitaries, and asks to show that this is necessarily the case. It also asks to show that
these measurements are optimal for noise resistance and for the Kullback--Leibler distance to local models; we do
not address those two parts here.

\subsection{What was known}
For $d=2$ the CGLMP expression is the CHSH expression and the statement is Tsirelson's bound~\cite{Tsirelson}.
For $d\ge 3$ the evidence was numerical~\cite{KGZMZ,DKZ,CGLMP}. For $d=3$, a semidefinite relaxation shows
that the maximal value on maximally entangled states of any dimension agrees with the DKZ value to within
$10^{-10}$~\cite{LVN} (this is the partial result listed in~\cite{OQP}). Ioannou and Rosset~\cite{IR} obtained
sum-of-squares certificates for the maximal quantum value of $I_3$ and $I_4$ over all states; by~\cite{ADGL} this
is a different optimisation problem from the one on maximally entangled states. In~\cite{repo} we gave exact
sum-of-squares certificates, with cyclotomic coefficients and checked by independent exact verifiers, which prove
that the DKZ measurements are optimal among projective measurements on maximally entangled states of any local
dimension for $d=3,\dots,20$, together with a rigidity statement. The
SATWAP expressions of Salavrakos et al.~\cite{SATWAP} are maximally violated by the maximally entangled state
with measurements of CGLMP type, admit sum-of-squares decompositions for all $d$, and underlie self-testing
results~\cite{KSTBSA,SSKA}; they are different Bell expressions and do not bear directly on $I_d$. Within the
family of Fourier measurements with four linear phase offsets, the DKZ point was recently found numerically to be
a strict local maximum for $d=2,\dots,6$~\cite[Sec.~V]{TMMD}. To our knowledge no argument valid for all $d$ was
available, even within a restricted class of measurements.

\subsection{The clock model and its commutative sector}
The starting point of this paper is an exact reduction, stated as Theorem~\ref{thm:reduction} below and proved
in Appendix~\ref{app:reduction} (it is also part of the material of~\cite{repo}). Put
\[
\psi_m=\frac{\pi m}{2d},\qquad h_m=\sec\psi_m-i\csc\psi_m\qquad(1\le m\le d-1),
\]
and for unitaries $V_0,\dots,V_{d-1}\in U(M)$ with $V_k^4=1$ let
\begin{equation}\label{eq:F}
F(V)=\sum_{0\le j<k\le d-1}\Re\bigl[h_{k-j}\,\tr(V_jV_k^*)\bigr],\qquad \tr=\tfrac1M\Tr .
\end{equation}
For every projective strategy on a maximally entangled state there is such a family with
$I_d=4F(V)/(d(d-1))$, and conversely every such family comes from a strategy. The DKZ strategy corresponds to
$V_0=\dots=V_{d-1}=1$, with
\begin{equation}\label{eq:FDKZ}
\FD=F(1,\dots,1)=\sum_{m=1}^{d-1}(d-m)\sec\psi_m,\qquad \IME(d)=\frac{4\FD}{d(d-1)} .
\end{equation}
So the optimality of the DKZ measurements on maximally entangled states (of any dimension, for projective
measurements) is equivalent to the inequality $F(V)\le\FD$ for all $M$ and all $d$-tuples of unitaries of order
dividing four.

When the $V_k$ commute, every joint eigenvector $v$ carries a configuration $a\in\Z_4^d$ with
$V_kv=i^{-a_k}v$, and $F(V)$ is an average of the values
\begin{equation}\label{eq:Fa}
F(a)=\sum_{0\le j<k\le d-1}\Re\bigl[h_{k-j}\,i^{a_k-a_j}\bigr],\qquad a\in\Z_4^d .
\end{equation}
This is a classical long-range four-state clock model on $d$ sites (equivalently, a long-range Ising ring with
antiperiodic boundary conditions, Section~\ref{sec:discussion}). The corresponding Bell strategies are genuinely
quantum: the DKZ strategy itself belongs to this commutative sector. We call $a$ \emph{one-step} if
$a_k=c+[k\ge r]$ for some $c\in\Z_4$ and $r\in\{0,\dots,d-1\}$; there are $4d$ one-step configurations, and
$a=0$ is one of them.

\subsection{Results}
For $N\ge2$ put $\kappa(u)=\cot(\pi u/N)$ for $u\in\Z_N\setminus\{0\}$ and $\kappa(0)=0$, and for $X,Y\subseteq\Z_N$
\[
\ip XY=\sum_{x\in X}\sum_{y\in Y}\kappa(x-y).
\]
For integers $p\le q$ let $[p,q)=\{p,p+1,\dots,q-1\}$, read modulo $N$. The function $\kappa$ is the kernel of the
discrete (periodic) Hilbert transform on $\Z_N$, so $\ip XY$ is its bilinear form on indicator functions.

\begin{introthm}[rearrangement inequality]\label{thm:A}
Let $N\ge 2$ and let $A,B\subseteq\Z_N$ be disjoint, $|A|=a$, $|B|=b$. Then
\[
\ip AB\le \Phi_N(a,b):=\ip{[0,a)}{[-b,0)} .
\]
If $a,b\ge1$ and $a+b<N$, equality holds if and only if $(A,B)=(r+[0,a),\,r+[-b,0))$ for some $r\in\Z_N$.
\end{introthm}

In words: the pairing is largest when $A$ and $B$ are intervals and $B$ ends exactly where $A$ begins. If $a=0$,
$b=0$ or $a+b=N$, both sides vanish (Section~\ref{sec:A}). The discrete Hilbert transform has mostly been studied
through Hilbert-type inequalities and operator norms (see, e.g.,~\cite{MV,BK}); we have not found
Theorem~\ref{thm:A}, or an equivalent statement, in the literature. Its continuous analogue
(Proposition~\ref{prop:cont}) is a direct consequence of known results~\cite{SW,Laeng} and the bathtub principle,
but the discrete statement does not follow from it (Remarks~\ref{rem:degenerate} and~\ref{rem:discrete}).

\begin{introthm}[commutative sector, all $d$]\label{thm:B}
Let $d\ge 2$ and let $V_0,\dots,V_{d-1}\in U(M)$ be pairwise commuting with $V_k^4=1$. Then $F(V)\le\FD$.
Equality holds if and only if every joint eigenvector $v$ of the $V_k$ satisfies $V_kv=i^{-a_k}v$ with a
one-step configuration $a$ (which may depend on $v$). The same holds for commuting unitaries of order dividing
four in any von Neumann algebra with a faithful normal tracial state $\tau$ in place of $\tr$.
\end{introthm}

For the physical statement, encode Alice's projective measurements $\{P^x_a\}_{a\in\Z_d}$ and Bob's
$\{Q^y_b\}_{b\in\Z_d}$ on $\C^D\otimes\C^D$ by the unitaries $U_x=\sum_aw^aP^x_a$ and $U'_y=\sum_bw^bQ^y_b$,
$w=e^{2\pi i/d}$, and put
\begin{equation}\label{eq:chainR}
R_1=U_2,\quad R_2=(U'_2)^T,\quad R_3=U_1,\quad R_4=(U'_1)^T,
\end{equation}
where the transpose is taken in the Schmidt basis of the maximally entangled state
$\Phi_D=D^{-1/2}\sum_{i=1}^D|ii\rangle$. The \emph{links} of the strategy are the following unitaries:
\[
L^{(0)}_n=R_1^nR_2^{-n},\quad L^{(1)}_n=R_2^nR_3^{-n},\quad L^{(2)}_n=R_3^nR_4^{-n},\quad
L^{(3)}_n=w^{-n}R_4^nR_1^{-n}\qquad(1\le n\le d-1);
\]
their normalised traces are the Fourier coefficients of the distributions of the outcome differences along the
chain $A_2,B_2,A_1,B_1$ (Section~\ref{sec:model}). We say that a strategy has \emph{equal links} if
$L^{(0)}_n=L^{(1)}_n=L^{(2)}_n=L^{(3)}_n$ for every $n$.

\begin{introcor}\label{cor:C}
Let $d\ge2$. Every projective strategy on a maximally entangled state $\Phi_D$ (any $D$) with equal links
satisfies $I_d\le\IME(d)$. Equality holds if and only if $d$ divides $D$ and, for some unitary
$u\colon\C^D\to\C^d\otimes\C^{D/d}$, the local unitary $u\otimes\bar u$ maps the strategy to the DKZ strategy
tensored with the identity on $\C^{D/d}$. In particular, each of the $4^d$ clock-phase strategies of
Section~\ref{subsec:sector} satisfies $I_d\le\IME(d)$, with equality exactly for the $4d$ one-step ones.
\end{introcor}

The equal-link strategies are exactly those whose reduced family $V$ in Theorem~\ref{thm:reduction} commutes
(Proposition~\ref{prop:links}). Theorem~\ref{thm:B} follows from Theorem~\ref{thm:A} through an exact identity
(Lemma~\ref{lem:cot}): with $N=4d$ and $S_a=\{k-da_k:0\le k\le d-1\}\subseteq\Z_{4d}$,
\[
F(a)=\tfrac12\ip{S_a}{S_a-d}-\tfrac d2 ,
\]
and $S_a$ is an interval exactly when $a$ is one-step. So Theorem~\ref{thm:B} is the case $N=4d$, $a=b=d$ of
Theorem~\ref{thm:A}.

\subsection{Method and difficulty}
The proof of Theorem~\ref{thm:A} (Section~\ref{sec:A}) has four steps. Let $C=\Z_N\setminus(A\cup B)$.
(i)~The three pairings $\ip AB,\ip BC,\ip CA$ are equal, and exchanging the colours of an adjacent pair of type
$(A,B)$, $(B,C)$ or $(C,A)$ increases the common value by an explicit positive amount. Hence every maximiser is
cyclically ordered: going around the circle the colours run through $w\ge1$ rounds of blocks $C,B,A$. The
configurations with $w=1$ are the rotations of $([0,a),[-b,0))$. (ii)~Smearing every point to a unit cell writes
$\ip AB$ as a continuous pairing plus a sum of explicit junction terms. (iii)~The continuous pairing is maximised
by adjacent arcs; this follows from the fact that the distribution of the conjugate function of the indicator
of a finite union of arcs, on the complement of the set, depends only on the total length (a circle version of
results of Stein and Weiss~\cite{SW} and Laeng~\cite{Laeng}, related to Boole's identity~\cite{Boole}), combined
with the bathtub principle. (iv)~A configuration with $w$ rounds has exactly $w$ junctions $B|A$, each of which costs about
$(2\log2-1)N/\pi$ in the junction terms; this beats all other junction contributions as soon as $w\ge2$.

Step~(iv) cannot be avoided. The continuous problem in step~(iii) has many maximisers, including the $w$-fold
covers of a pair of adjacent arcs (Remark~\ref{rem:degenerate}), so the discrete inequality is decided by the
discretisation error, which is of order $N$. In the language of the clock model this is a near-degeneracy between
the DKZ configuration and ``$w$-fold'' configurations (Section~\ref{sec:discussion}), and it suggests why
certificates that are local in the site index fail for large $d$: for instance, the Sherali--Adams relaxation
of the clock model built from all three-site marginals is exact for $d\le8$ but not for $9\le d\le12$
(numerical; Section~\ref{sec:discussion}).

\subsection{What remains open}
Theorem~\ref{thm:B} covers commuting families, and Theorem~\ref{thm:D} in Section~\ref{sec:discussion} covers the
non-commuting two-window families. Whether $F(V)\le\FD$ holds for all non-commuting families, which is what remains
of the optimality statement of Problem~27B on maximally entangled states, is open for $d\ge21$. In
Section~\ref{sec:discussion} we formulate an operator version of Theorem~\ref{thm:A} that would imply it. The
cyclic symmetry extends to operators verbatim and the exchange identity only partially; we do not know operator
versions of the continuous inequality and of the junction estimate.

\subsection{Organisation}
Section~\ref{sec:model} recalls the reduction to the clock model and characterises the commutative sector.
Section~\ref{sec:cot} proves the cotangent representation and deduces Theorem~\ref{thm:B} and
Corollary~\ref{cor:C} from Theorem~\ref{thm:A}. Section~\ref{sec:A} proves Theorem~\ref{thm:A}.
Section~\ref{sec:numerics} lists numerical cross-checks, which are not part of the proofs, and
Section~\ref{sec:discussion} discusses the non-commutative case. Appendix~\ref{app:reduction} contains the proof
of the reduction.

\section{The CGLMP expression and the clock model}\label{sec:model}

Throughout, $d\ge2$, $w=e^{2\pi i/d}$, $z=e^{2\pi i/(4d)}$ (so $z^4=w$ and $z^d=i$), and $m(t)\in\{0,\dots,d-1\}$
denotes the residue of an integer $t$ modulo $d$.

\subsection{A chain form of the CGLMP expression}
Relabel the settings as $X_1=A_2$, $Y_1=B_2$, $X_2=A_1$, $Y_2=B_1$, and put
\begin{equation}\label{eq:S}
S=\mathbb{E}\,m(X_1-Y_1)+\mathbb{E}\,m(Y_1-X_2)+\mathbb{E}\,m(X_2-Y_2)+\mathbb{E}\,m(Y_2-X_1-1),
\end{equation}
each expectation being taken in the distribution of the pair of outcomes that occurs in it.

\begin{lemma}\label{lem:chain}
For every behaviour $P(a,b|x,y)$, $I_d=4-\frac{2}{d-1}\,S$.
\end{lemma}

\begin{proof}
Consider the terms of \eqref{eq:cglmp} that involve $A_1$ and $B_1$, and let $t=m(A_1-B_1)$. For
$0\le k\le\lfloor d/2\rfloor-1$ the event $t=k$ has weight $1-\frac{2k}{d-1}$, and the event $t=d-1-k$ has weight
$-(1-\frac{2k}{d-1})=1-\frac{2(d-1-k)}{d-1}$. If $d$ is odd, the remaining value $t=(d-1)/2$ does not occur and
$1-\frac{2t}{d-1}=0$ there. So these terms equal $\mathbb{E}f(t)$ with $f(t)=1-\frac{2t}{d-1}$ for all
$t\in\{0,\dots,d-1\}$, that is $1-\frac{2}{d-1}\mathbb{E}\,m(A_1-B_1)$. In the same way the terms with $(B_1,A_2)$,
$(A_2,B_2)$ and $(B_2,A_1)$ give $1-\frac2{d-1}\mathbb{E}\,m(B_1-A_2-1)$, $1-\frac2{d-1}\mathbb{E}\,m(A_2-B_2)$ and
$1-\frac2{d-1}\mathbb{E}\,m(B_2-A_1)$. In the relabelled variables the sum of the four expectations is $S$.
\end{proof}

\begin{remark}
The local bound $I_d\le2$ reads $S\ge d-1$; this is the form in which the inequality is written in~\cite{OQP},
where it is attributed to Gill. For no-signalling behaviours one also has $S=d\Pi-1$ with
$\Pi=P(X_1<Y_1)+P(Y_1<X_2)+P(X_2<Y_2)+P(Y_2\le X_1)$ (outcomes compared as integers in $\{0,\dots,d-1\}$), which
follows from $m(u-v)=u-v+d\,[u<v]$ for $u,v\in\{0,\dots,d-1\}$. So $I_d\le2$ is equivalent to $\Pi\ge1$, the form
of the CGLMP inequality used by Zohren and Gill~\cite[eq.~(3)]{ZG}. We do not use these facts.
\end{remark}

\subsection{Maximally entangled strategies}
A projective strategy on $\Phi_D$ gives $P(a,b|x,y)=\langle\Phi_D|P^x_a\otimes Q^y_b|\Phi_D\rangle
=\tau\bigl(P^x_a(Q^y_b)^T\bigr)$ with $\tau=\frac1D\Tr$. With the unitaries \eqref{eq:chainR}, which satisfy
$R_i^d=1$, and the Fourier expansion
\begin{equation}\label{eq:mfourier}
m(t)=\frac{d-1}{2}+\sum_{n=1}^{d-1}c_n\,w^{nt},\qquad c_n=-\frac{1}{1-w^{-n}},
\end{equation}
one finds $\mathbb{E}\,w^{n(X_1-Y_1)}=\tau(R_1^nR_2^{-n})$, $\mathbb{E}\,w^{n(Y_1-X_2)}=\tau(R_2^nR_3^{-n})$,
$\mathbb{E}\,w^{n(X_2-Y_2)}=\tau(R_3^nR_4^{-n})$ and $\mathbb{E}\,w^{n(Y_2-X_1-1)}=w^{-n}\tau(R_4^nR_1^{-n})$.
Hence
\begin{equation}\label{eq:Strace}
S=2(d-1)+\sum_{n=1}^{d-1}c_n\,\Lambda_n(R),\qquad
\Lambda_n(R)=\sum_{i=0}^{3}\tau\bigl(L^{(i)}_n\bigr),
\end{equation}
with the links $L^{(i)}_n$ defined in the introduction. Conversely, any four unitaries $R_1,\dots,R_4$ on $\C^D$
with $R_i^d=1$ come from a projective strategy on $\Phi_D$, through their spectral projections. (The coefficient
in \eqref{eq:mfourier} is $\frac1d\sum_{t=0}^{d-1}t\,w^{-nt}=-1/(1-w^{-n})$ for $n\ne0$.)

\subsection{The reduction}
For unitaries $V_0,\dots,V_{d-1}\in U(M)$ with $V_k^4=1$ let $X$ be the cyclic shift $X|k\rangle=|k+1\rangle$ on
$\C^d$ (indices modulo $d$) and define the \emph{covariant strategy} $R^V$ on $\C^d\otimes\C^M$ by
\begin{equation}\label{eq:covariant}
W=\sum_{k=0}^{d-1}|k\rangle\langle k|\otimes z^kV_k,\qquad R^V_1=X\otimes\one_M,\qquad
R^V_{i+1}=WR^V_iW^*\quad(i=1,2,3).
\end{equation}

\begin{theorem}[reduction to the clock model]\label{thm:reduction}
\begin{enumerate}
\item For every covariant strategy, $S(R^V)=2(d-1)-\frac2d F(V)$, that is $I_d(R^V)=\frac{4F(V)}{d(d-1)}$.
\item For every projective strategy $R$ on $\Phi_D$ there are $M$ (one can take $M=4D$) and unitaries
$V_0,\dots,V_{d-1}\in U(M)$ with $V_k^4=1$ such that $R$ is unitarily equivalent to a direct summand of
$R^V$ (simultaneous conjugation of $R_1,\dots,R_4$) and $I_d(R)=I_d(R^V)=\frac{4F(V)}{d(d-1)}$.
\end{enumerate}
Consequently, $I_d\le\IME(d)$ holds for all projective strategies on maximally entangled states of all
dimensions if and only if $F(V)\le\FD$ for all $M$ and all $V_0,\dots,V_{d-1}\in U(M)$ with $V_k^4=1$.
\end{theorem}

The family $V$ in (2) is produced by a twirl over a cyclic symmetry of order $4d$ followed by the
Stone--von Neumann theorem for $\Z_d$; it is unique up to simultaneous unitary conjugation of the $V_k$
(Appendix~\ref{app:reduction}). We call it the \emph{reduced family} of $R$.

\subsection{The DKZ strategy}\label{subsec:dkz}
Let $W_0=\mathrm{diag}(z^k)_{k=0}^{d-1}$, so that $R^{(0)}:=R^V$ for $M=1$, $V\equiv1$, is
$R^{(0)}_i=W_0^{i-1}XW_0^{-(i-1)}$. The measurement procedure of~\cite[eqs.~(12)--(15)]{CGLMP} (diagonal phases
$e^{2\pi i\alpha_xj/d}$ for Alice and $e^{2\pi i\beta_yj/d}$ for Bob, then a discrete Fourier transform on each
side, then a measurement in the computational basis, with $\alpha=(0,\frac12)$ and $\beta=(\frac14,-\frac14)$)
amounts to measuring in the orthonormal bases
$|k\rangle_{A,x}=d^{-1/2}\sum_je^{-2\pi ij(k+\alpha_x)/d}|j\rangle$ and
$|l\rangle_{B,y}=d^{-1/2}\sum_je^{2\pi ij(l-\beta_y)/d}|j\rangle$; this gives the joint probabilities
$P(A_x=k,B_y=l)=\bigl(2d^3\sin^2(\pi(k-l+\alpha_x+\beta_y)/d)\bigr)^{-1}$ of~\cite[eq.~(15)]{CGLMP}. Writing
$D_\gamma=\mathrm{diag}(w^{\gamma j})$ and $F|j\rangle=d^{-1/2}\sum_kw^{jk}|k\rangle$ one has
$U_x=D_{\alpha_x}^*(F^*Z_0F)D_{\alpha_x}=D_{\alpha_x}^*XD_{\alpha_x}$ with $Z_0=\mathrm{diag}(w^k)$, and similarly
$(U'_y)^T=D_{\beta_y}XD_{\beta_y}^*$. Since $D_{1/4}=W_0$, the chain \eqref{eq:chainR} of the DKZ strategy is
\[
(R_1,R_2,R_3,R_4)=(W_0^{-2}XW_0^{2},\;W_0^{-1}XW_0,\;X,\;W_0XW_0^{-1})=G\,R^{(0)}\,G^*,\qquad G=W_0^{-2}.
\]
Thus DKZ is, up to the local unitary $G\otimes\bar G$ (which fixes $\Phi_d$), the covariant strategy with
$V\equiv1$, and $I_d(\mathrm{DKZ})=4\FD/(d(d-1))=\IME(d)$ by Theorem~\ref{thm:reduction}(1).

\subsection{The commutative sector}\label{subsec:sector}
For $a\in\Z_4^d$ the \emph{clock-phase strategy} $R^a$ on $\Phi_d$ is the covariant strategy with $M=1$ and
$V_k=i^{-a_k}$, that is
\[
R^a_1=X,\qquad R^a_{i+1}=W_aR^a_iW_a^*,\qquad W_a=\mathrm{diag}\bigl(z^k\,i^{-a_k}\bigr)_{k=0}^{d-1}.
\]
All four of its measurement bases are Fourier bases dressed by diagonal phases, consecutive ones differing by the
same diagonal unitary $W_a$, and $I_d(R^a)=4F(a)/(d(d-1))$ with $F(a)$ as in \eqref{eq:Fa}. If $V_0,\dots,V_{d-1}$
commute, a common eigenbasis $(v_\mu)$ splits $R^V$ into the direct sum $\bigoplus_\mu R^{a(\mu)}$, where
$V_kv_\mu=i^{-a_k(\mu)}v_\mu$.

\begin{proposition}[equal links]\label{prop:links}
A projective strategy on $\Phi_D$ has equal links if and only if its reduced family consists of pairwise
commuting unitaries. In particular every clock-phase strategy, and every covariant strategy with commuting $V$,
has equal links.
\end{proposition}

The proof is in Appendix~\ref{app:reduction}. The equal-link condition for $n=1$ is one of the conditions that
appear in the rigidity analysis of~\cite{repo}.

\section{The cotangent representation; proofs of Theorem~\ref{thm:B} and Corollary~\ref{cor:C}}\label{sec:cot}

\begin{lemma}[cotangent representation]\label{lem:cot}
Let $N=4d$ and $\kappa(u)=\cot(\pi u/N)$ as above. For $a\in\Z_4^d$ put
$S_a=\{k-da_k \bmod N: 0\le k\le d-1\}$. Then $S_a$ contains exactly one element of each residue class modulo
$d$, every such transversal is $S_a$ for exactly one $a$, the sets $S_a$ and $S_a-d$ are disjoint, and
\[
F(a)=\frac12\sum_{\substack{x,y\in S_a\\ x\ne y}}\sec\frac{\pi(x-y)}{2d}=\frac12\ip{S_a}{S_a-d}-\frac d2 .
\]
Moreover $S_a$ is an interval $[c,c+d)$ if and only if $a$ is one-step, and then $(S_a,S_a-d)=([c,c+d),[c-d,c))$.
\end{lemma}

\begin{proof}
Let $x_k=k-da_k$. For $j<k$, $m=k-j$ and $s=a_k-a_j$ we have $x_k-x_j=m-ds$ and
$\Re[h_mi^s]=\sec(\psi_m-s\pi/2)$ (for $s=0,1,2,3$ both sides equal $\sec\psi_m,\csc\psi_m,-\sec\psi_m,-\csc\psi_m$),
which is $\sec(\pi(x_k-x_j)/(2d))$; the function $u\mapsto\sec(\pi u/(2d))$ is even and has period $4d$. This
gives the first expression for $F(a)$. The first assertions are clear since $x_k\equiv k\pmod d$ and $a_k$ ranges
over $\Z_4$.

Next, $\cot p-\cot q=\sin(q-p)/(\sin p\sin q)$ and $\sin(\frac t2+\frac\pi4)\sin(\frac t2-\frac\pi4)=-\frac12\cos t$
give $\sec t=\frac12[\cot(\frac t2+\frac\pi4)-\cot(\frac t2-\frac\pi4)]$. With $t=2\pi(x-y)/N$ and $\pi/4=\pi d/N$
this reads
\[
\sec\frac{\pi(x-y)}{2d}=\frac12\bigl[\kappa(x-y+d)-\kappa(x-y-d)\bigr]\qquad(x\not\equiv y\bmod d),
\]
where no pole occurs because $x-y\pm d\not\equiv0\pmod{4d}$. Summing over ordered pairs $x\ne y$ in $S_a$ and using
that $\kappa$ is odd, the two halves are equal, so
$F(a)=\frac12\sum_{x\ne y}\kappa(x-y+d)=\frac12\bigl[\ip{S_a}{S_a-d}-d\,\kappa(d)\bigr]$ and $\kappa(d)=\cot(\pi/4)=1$.
If $x=y-d$ with $x,y\in S_a$, then $x\equiv y\pmod d$, so $x=y$ and $d\equiv0\pmod{4d}$, which is impossible; hence
$S_a\cap(S_a-d)=\emptyset$. Finally, if $a_k=c+[k\ge r]$ then
$S_a=\{-dc+k:k<r\}\cup\{-dc-d+k:r\le k<d\}=-dc+[r-d,r)$. Conversely an interval of length $d$ is a transversal, so it
is $S_a$ for a unique $a$; the $4d$ intervals and the $4d$ one-step configurations therefore correspond to each
other bijectively.
\end{proof}

\begin{proof}[Proof of Theorem~\ref{thm:B} from Theorem~\ref{thm:A}]
Let $a\in\Z_4^d$. By Lemma~\ref{lem:cot} and Theorem~\ref{thm:A} with $N=4d$, $A=S_a$, $B=S_a-d$ (so that
$|A|=|B|=d$ and $|C|=2d\ge1$),
\[
F(a)=\tfrac12\ip{S_a}{S_a-d}-\tfrac d2\le\tfrac12\Phi_{4d}(d,d)-\tfrac d2=F(0)=\FD ,
\]
because $S_0=[0,d)$ and $S_0-d=[-d,0)$. Equality holds if and only if $(S_a,S_a-d)$ is a rotation of $([0,d),[-d,0))$,
that is, if and only if $S_a$ is an interval, that is, if and only if $a$ is one-step.

Now let $V_0,\dots,V_{d-1}$ commute, $V_k^4=1$, in a von Neumann algebra with faithful normal tracial state $\tau$
(for matrices, $\tau=\tr$). The spectral projections $P_{k,s}$ of $V_k$ ($V_k=\sum_{s\in\Z_4}i^{-s}P_{k,s}$) commute,
so $\Pi_a=\prod_kP_{k,a_k}$, $a\in\Z_4^d$, is a finite family of orthogonal projections with sum $1$ and
$V_k=\sum_ai^{-a_k}\Pi_a$. Then $\tau(V_jV_k^*)=\sum_a\tau(\Pi_a)\,i^{a_k-a_j}$ and
\[
F(V)=\sum_{a\in\Z_4^d}\tau(\Pi_a)\,F(a)\le\FD ,
\]
since $\tau(\Pi_a)\ge0$ and $\sum_a\tau(\Pi_a)=1$. Equality forces $\tau(\Pi_a)=0$, hence $\Pi_a=0$, for every $a$
that is not one-step. For matrices this is the statement about joint eigenvectors.
\end{proof}

\begin{proof}[Proof of Corollary~\ref{cor:C}]
Let $R$ have equal links and let $V$ be its reduced family. By Proposition~\ref{prop:links} the $V_k$ commute,
so by Theorems~\ref{thm:reduction} and~\ref{thm:B}, $I_d(R)=4F(V)/(d(d-1))\le4\FD/(d(d-1))=\IME(d)$.

Suppose equality holds. Diagonalise the $V_k$ in a common orthonormal basis $(v_\mu)_{\mu=1}^M$ (this conjugates
$W$ by $\one\otimes u_0$ and leaves $X\otimes\one$ unchanged). By Theorem~\ref{thm:B} every $a(\mu)$ is one-step,
$a_k(\mu)=c_\mu+[k\ge r_\mu]$, and $R^V=\bigoplus_\mu R^{a(\mu)}$. For $0\le r\le d-1$ one computes
$X^rW_0X^{-r}=\mathrm{diag}(z^{(k-r)\bmod d})=z^{-r}i\,\mathrm{diag}(z^ki^{-[k\ge r]})$, using $z^d=i$. Hence
$W_{a(\mu)}=\gamma_\mu X^{r_\mu}W_0X^{-r_\mu}$ with $|\gamma_\mu|=1$, and since $X$ commutes with $X^{r_\mu}$,
$R^{a(\mu)}_i=X^{r_\mu}R^{(0)}_iX^{-r_\mu}$ for $i=1,\dots,4$. So $R^V$ is unitarily equivalent to
$R^{(0)}\otimes\one_M$, and by Theorem~\ref{thm:reduction}(2) $R$ is unitarily equivalent to the restriction of
$R^{(0)}\otimes\one_M$ to an invariant subspace $H'$. The $*$-algebra generated by $R^{(0)}_1,\dots,R^{(0)}_4$ is all of
$M_d(\C)$: it contains $X$ and $R^{(0)}_1(R^{(0)}_2)^{-1}=XW_0X^{-1}W_0^*=z^{-1}(\one+(i-1)|0\rangle\langle0|)$,
hence $|0\rangle\langle0|$ and all matrix units $X^k|0\rangle\langle0|X^{-l}$. Therefore the commutant of
$\{R^{(0)}_i\otimes\one_M\}$ is $\one\otimes M_M(\C)$, and the invariant subspace $H'$, being the range of a
projection in this commutant, is $\C^d\otimes K$ for a subspace $K$. So $D=d\dim K$ and $R$ is unitarily
equivalent to $R^{(0)}\otimes\one_K$, hence (Section~\ref{subsec:dkz}) to $\mathrm{DKZ}\otimes\one_K$.
A simultaneous conjugation $R_i\mapsto uR_iu^*$ of the unitaries \eqref{eq:chainR} is the same as the local unitary
$u\otimes\bar u$ applied to the strategy, and $(u\otimes\bar u)\Phi_D$ is again maximally entangled. Conversely
$\mathrm{DKZ}\otimes\one$ attains $\IME(d)$ and has equal links.

For the clock-phase strategies, $I_d(R^a)=4F(a)/(d(d-1))$, and $F(a)\le\FD$ with equality exactly for one-step $a$
by the first part of the proof of Theorem~\ref{thm:B}.
\end{proof}

\section{Proof of Theorem~\ref{thm:A}}\label{sec:A}

Fix $N\ge2$, write $\kappa(u)=\cot(\pi u/N)$ also for real $u\notin N\Z$, and recall $\kappa(0)=0$ on $\Z_N$. The
function $\kappa$ is odd and $N$-periodic. For $X\subseteq\Z_N$ let
\[
p_X(t)=\sum_{s\in X}\kappa(t-s)\qquad(t\in\Z_N).
\]
Two facts will be used repeatedly: $\ip YX=-\ip XY$ (oddness; in particular $\ip XX=0$), and
$\sum_{u\in\Z_N}\kappa(u)=0$ (the terms $u$ and $-u$ cancel, and $\kappa(N/2)=0$ if $N$ is even), so that
$\ip X{\Z_N}=0$ for every $X$.

Given disjoint $A,B$ we always put $C=\Z_N\setminus(A\cup B)$ and $c=|C|$, and think of the elements of $A$, $B$,
$C$ as coloured by the three colours $A$, $B$, $C$.

\emph{Trivial cases.} If $a=0$ or $b=0$, both $\ip AB$ and $\Phi_N(a,b)$ are empty sums. If $c=0$, then
$\ip AB=\ip A{\Z_N}-\ip AA=0$, and $\Phi_N(a,b)=0$ for the same reason. So from now on $a,b,c\ge1$, and in
particular $N\ge3$.

\subsection{Cyclic symmetry}

\begin{lemma}[cyclic symmetry]\label{lem:cyclic}
The three pairings are equal, $\ip AB=\ip BC=\ip CA=:Q$. Consequently
\[
Q=\sum_{y\in B}p_C(y)=\sum_{t\in C}p_A(t)=\sum_{x\in A}p_B(x).
\]
\end{lemma}

\begin{proof}
$0=\ip A{\Z_N}=\ip AA+\ip AB+\ip AC$ gives $\ip AB=-\ip AC=\ip CA$, and $0=\ip B{\Z_N}$ gives
$\ip BC=-\ip BA=\ip AB$. The three formulas are $\ip BC$, $\ip CA$ and $\ip AB$ written out.
\end{proof}

\subsection{The exchange identity}
For $1\le u\le N-2$ let $\delta(u)=\kappa(u)-\kappa(u+1)$. Since $\pi u/N$ and $\pi(u+1)/N$ lie in $(0,\pi)$ and
$\cot$ is strictly decreasing there, $\delta(u)>0$.

\begin{lemma}[exchange identity]\label{lem:swap}
Let $x\in\Z_N$ with $x\in A$ and $x+1\in B$, and let $(A',B')$ be obtained by exchanging the two colours
($A'=A\setminus\{x\}\cup\{x+1\}$, $B'=B\setminus\{x+1\}\cup\{x\}$). Then
\[
\ip{A'}{B'}-\ip AB=p_C(x)-p_C(x+1)=\sum_{t\in C}\delta(x-t)>0,
\]
where $x-t$ is read in $\{1,\dots,N-2\}$. Likewise, exchanging an adjacent pair coloured $(B,C)$ increases the
common value $Q$ of Lemma~\ref{lem:cyclic} by $\sum_{s\in A}\delta(x-s)>0$, and exchanging an adjacent pair
coloured $(C,A)$ increases it by $\sum_{y\in B}\delta(x-y)>0$.
\end{lemma}

\begin{proof}
The exchange does not change $C$, so by Lemma~\ref{lem:cyclic} (applied before and after)
$\ip{A'}{B'}-\ip AB=\sum_{y\in B'}p_C(y)-\sum_{y\in B}p_C(y)=p_C(x)-p_C(x+1)$. For $t\in C$ we have
$t\notin\{x,x+1\}$, so $x-t\bmod N\in\{1,\dots,N-2\}$ and $\kappa(x-t)-\kappa(x+1-t)=\delta(x-t)$. The sum is
positive because $C\neq\emptyset$. For a pair coloured $(B,C)$ the set $A$ is unchanged and
$Q=\sum_{t\in C}p_A(t)$; for a pair coloured $(C,A)$ the set $B$ is unchanged and $Q=\sum_{x\in A}p_B(x)$. The same
computation applies, and $A,B\ne\emptyset$.
\end{proof}

We call a colouring \emph{cyclically ordered} if no adjacent pair $(x,x+1)$ is coloured $(A,B)$, $(B,C)$ or $(C,A)$.

\begin{corollary}\label{cor:monotone}
Let $a,b,c\ge1$. Every maximiser of $\ip AB$ (over disjoint $A,B$ of sizes $a,b$) is cyclically ordered. A
cyclically ordered colouring consists of $3w$ nonempty maximal blocks of consecutive points, for some integer
$w\ge1$ (the winding number), which occur in the cyclic order
\[
C_1,B_1,A_1,\;C_2,B_2,A_2,\;\dots,\;C_w,B_w,A_w
\]
in the direction of increasing $x$. The colourings with $w=1$ are exactly the rotations of $([0,a),[-b,0))$.
\end{corollary}

\begin{proof}
The first assertion follows from Lemma~\ref{lem:swap}, since exchanges preserve the sizes. In a cyclically
ordered colouring every colour change from $x$ to $x+1$ is of type $B\to A$, $A\to C$ or $C\to B$; assigning
$C\mapsto0$, $B\mapsto1$, $A\mapsto2$ in $\Z_3$, each change adds $1$. Going once around the circle the total
change is $0$ in $\Z_3$, so the number of changes is $3w$ for an integer $w$, and $w\ge1$ because all three colours
occur. The maximal blocks between consecutive changes then follow the stated order. For $w=1$ there is one block
of each colour, in the order $C,B,A$; so $B=r+[-b,0)$ and $A=r+[0,a)$ for some $r$.
\end{proof}

\subsection{Smearing}
Let $\T_N=\R/N\Z$ with Lebesgue measure. For $X\subseteq\Z_N$ let $\widetilde X=\bigcup_{x\in X}[x,x+1)\subseteq\T_N$,
and for measurable $U,V\subseteq\T_N$ put
\[
\Qc(U,V)=\int_U\int_V\kappa(\xi-\eta)\,d\eta\,d\xi
\]
whenever the integral converges absolutely. For $m\in\{1,\dots,N-1\}$ define
\[
\kb(m)=\int_0^1\!\!\int_0^1\kappa(m+s-t)\,ds\,dt=\int_{-1}^{1}(1-|v|)\,\kappa(m+v)\,dv,\qquad e(m)=\kappa(m)-\kb(m).
\]
The integrand of $\kb(m)$ is bounded: the only possible poles, at $v=-1$ for $m=1$ and at $v=1$ for $m=N-1$, sit
where the weight $1-|v|$ vanishes to first order.

\begin{lemma}[smearing]\label{lem:smear}
For disjoint $A,B\subseteq\Z_N$ the integral $\Qc(\widetilde A,\widetilde B)$ converges absolutely and
\[
\ip AB=\Qc(\widetilde A,\widetilde B)+E(A,B),\qquad E(A,B)=\sum_{x\in A,\,y\in B}e(x-y),
\]
where $x-y$ is read in $\{1,\dots,N-1\}$.
\end{lemma}

\begin{proof}
$\Qc(\widetilde A,\widetilde B)=\sum_{x\in A,y\in B}\int_0^1\int_0^1\kappa(x-y+s-t)\,ds\,dt$. For $x-y\equiv m$ the
integrand is singular only if $m+s-t\in N\Z$, which for $s,t\in[0,1]$ happens only at a corner of the square
($m=1$, $s=0$, $t=1$, or $m=N-1$, $s=1$, $t=0$). There $|\kappa|$ is bounded by a constant times
$1/(s+1-t)$, respectively $1/(1-s+t)$, which is integrable on the square. The substitution $v=s-t$, whose density
on $[-1,1]$ is $1-|v|$, gives the second expression for $\kb$, and the cell integral equals $\kb(x-y)$.
\end{proof}

\subsection{The continuous inequality}
Let $\T=\R/2\pi\Z$. A finite union $V'$ of pairwise disjoint open arcs $(p_j,q_j)$, $j=1,\dots,n$, with
$p_j<q_j<p_j+2\pi$ and total length $\beta=\sum_j(q_j-p_j)\in(0,2\pi)$, has the conjugate function
\begin{equation}\label{eq:conj}
g_{V'}(\theta)=\frac1{2\pi}\,\mathrm{p.v.}\!\int_{V'}\cot\frac{\theta-\varphi}{2}\,d\varphi
=\frac1\pi\sum_{j=1}^n\log\left|\frac{\sin\frac{\theta-p_j}2}{\sin\frac{\theta-q_j}2}\right|
\qquad(\theta\notin\{p_j,q_j\}),
\end{equation}
since $\frac{d}{d\varphi}\bigl(-2\log|\sin\frac{\theta-\varphi}2|\bigr)=\cot\frac{\theta-\varphi}2$. It has logarithmic
singularities at the endpoints and is integrable on $\T$.

\begin{lemma}[distribution of the conjugate function]\label{lem:boole}
In this situation, for every $\lambda\in\R$,
\[
\bigl|\{\theta\in\T\setminus V':g_{V'}(\theta)>\lambda\}\bigr|=2\pi\,\mu_\beta\bigl((0,e^{-\pi\lambda})\bigr),
\]
where $\mu_\beta$ is the Cauchy distribution on $\R$ with density
$\frac1\pi\,\frac{\sin(\beta/2)}{(t-\cos(\beta/2))^2+\sin^2(\beta/2)}$. In particular the left side depends only on
$\beta$ and $\lambda$.
\end{lemma}

\begin{proof}
Let $\D$ be the open unit disc and $\HH$ the open upper half plane. Consider the rational function
\[
\Psi(z)=e^{-i\beta/2}\prod_{j=1}^n\frac{e^{iq_j}-z}{e^{ip_j}-z}.
\]
\emph{$\Psi$ maps $\D$ into $\HH$.} Let $\mathcal F(z)=\frac1{2\pi}\int_{V'}\frac{e^{i\varphi}+z}{e^{i\varphi}-z}\,d\varphi$
($z\in\D$), an analytic function whose real part is the Poisson integral of $\one_{V'}$. As $0<|V'|<2\pi$, the
positivity of the Poisson kernel gives $0<\Re\mathcal F<1$ on $\D$. For one arc $(p,q)$, the partial fractions
$\frac{\zeta+z}{\zeta(\zeta-z)}=\frac2{\zeta-z}-\frac1\zeta$ and the substitution $\zeta=e^{i\varphi}$ give
$\frac1{2\pi}\int_{p}^{q}\frac{e^{i\varphi}+z}{e^{i\varphi}-z}d\varphi
=\frac1{\pi i}\log\frac{e^{iq}-z}{e^{ip}-z}-\frac{q-p}{2\pi}$ for some branch of the logarithm.
Summing over the arcs and exponentiating, $\exp(i\pi\mathcal F)=\Psi$ on $\D$, whatever the branches. Hence
$\arg\Psi=\pi\Re\mathcal F\in(0,\pi)$, that is, $\Psi(\D)\subseteq\HH$. Also $\Psi(0)=e^{-i\beta/2}e^{i\beta}=e^{i\beta/2}$.

\emph{Boundary values.} For $\theta\notin\{p_j,q_j\}$,
$e^{iq_j}-e^{i\theta}=-2ie^{i(q_j+\theta)/2}\sin\frac{\theta-q_j}2$ and similarly for $p_j$, so
\[
\Psi(e^{i\theta})=\prod_{j=1}^n\frac{\sin\frac{\theta-q_j}2}{\sin\frac{\theta-p_j}2}\in\R\setminus\{0\},\qquad
|\Psi(e^{i\theta})|=e^{-\pi g_{V'}(\theta)} .
\]
Each factor is $2\pi$-periodic in $\theta$; taking $\theta\in(p_j,p_j+2\pi)$, its denominator is positive and its
numerator is negative exactly for $\theta<q_j$. So the $j$-th factor is negative exactly on the $j$-th arc,
and, the arcs being disjoint, $\Psi(e^{i\theta})<0$ if and only if $\theta\in V'$. Therefore
\[
\{\theta\in\T\setminus V':g_{V'}(\theta)>\lambda\}=\{\theta:0<\Psi(e^{i\theta})<e^{-\pi\lambda}\}
\]
up to finitely many points.

\emph{The law of the boundary values.} Let $w_0=\Psi(0)=e^{i\beta/2}\in\HH$ and
$\mathcal M(\zeta)=(\zeta-w_0)/(\zeta-\bar w_0)$, a M\"obius map of $\HH$ onto $\D$ and of $\R\cup\{\infty\}$
onto the unit circle, with $\mathcal M(w_0)=0$. The rational function $B=\mathcal M\circ\Psi$ maps $\D$ into $\D$;
it has no poles on the closed disc (it could only have poles where $\Psi=\bar w_0$, which does not happen there;
at the poles $e^{ip_j}$ of $\Psi$ it takes the value $1$), and $|B|=1$ on the unit circle. For $k\ge1$ the mean
value property of $B^k$ gives $\frac1{2\pi}\int_0^{2\pi}B(e^{i\theta})^k\,d\theta=B(0)^k=0$, and by conjugation also
$\frac1{2\pi}\int\overline{B}^k=0$. So the image of the normalised arc length under $\theta\mapsto B(e^{i\theta})$ has
the Fourier coefficients of the normalised arc length, and is therefore the normalised arc length. Finally,
$\Psi(e^{i\theta})=\mathcal M^{-1}(B(e^{i\theta}))$ off finitely many points, and $\mathcal M^{-1}$ carries the
normalised arc length to the harmonic measure of $\HH$ at $\mathcal M^{-1}(0)=w_0$, which is the Cauchy distribution
with centre $\Re w_0=\cos(\beta/2)$ and scale $\Im w_0=\sin(\beta/2)$. (For a bounded continuous $f$ on $\R$ with a
limit at $\infty$, its Poisson extension $u$ to $\HH$ satisfies $u(w_0)=\frac1{2\pi}\int f(\mathcal M^{-1}(e^{it}))\,dt$ by
the mean value property of the harmonic function $u\circ\mathcal M^{-1}$ at $0$.) So the image of $d\theta/2\pi$ under
$\theta\mapsto\Psi(e^{i\theta})$ is $\mu_\beta$, and the lemma follows.
\end{proof}

For sets $E$ of finite measure on the line, the fact that the distribution function of the Hilbert transform of
$\one_E$ depends only on $|E|$ is a theorem of Stein and Weiss~\cite{SW}; for sums of point masses the analogous
statement is Boole's identity~\cite{Boole} (see also Loomis~\cite{Loomis} and~\cite{CLM}). That the distributions
on $E$ and on its complement depend only on $|E|$ separately is due to Laeng~\cite{Laeng} (see the discussion
in~\cite[Section~4]{GL}). Lemma~\ref{lem:boole} is the circle version of the latter statement for finite unions of
arcs, which is all we need; the proof above is the standard one, made explicit by the fact that $\Psi$ is a
rational function. Proposition~\ref{prop:cont} below is then a direct consequence of the bathtub principle, and
we do not claim novelty for it.

\begin{proposition}[continuous rearrangement]\label{prop:cont}
Let $U,V\subseteq\T_N$ be disjoint, each a finite union of half-open intervals, with $|U|=a$ and $|V|=b$, where
$0<b<N$. Then $\Qc(U,V)$ converges absolutely and
\[
\Qc(U,V)\le\Qc\bigl([0,a),[-b,0)\bigr).
\]
\end{proposition}

\begin{proof}
Absolute convergence: the closures of $U$ and $V$ meet in finitely many points. Near such a point $|\kappa(\xi-\eta)|$
is at most a constant times $1/|\xi-\eta|$, with $\xi$ and $\eta$ on opposite sides of the point, and
$\int_0^{1}\int_0^{1}\frac{ds\,dt}{s+t}<\infty$; elsewhere the integrand is bounded.

Rescale by $\theta=2\pi\xi/N$, $\varphi=2\pi\eta/N$, and let $U',V'\subseteq\T$ be the images of $U,V$; after merging
touching intervals and discarding finitely many points, $V'$ is a union of disjoint open arcs of total length
$\beta=2\pi b/N\in(0,2\pi)$, and $U'\subseteq\T\setminus V'$ has measure $\alpha=2\pi a/N$. By Fubini and
\eqref{eq:conj},
\[
\Qc(U,V)=\Bigl(\frac N{2\pi}\Bigr)^2\int_{U'}\int_{V'}\cot\frac{\theta-\varphi}2\,d\varphi\,d\theta
=\frac{N^2}{2\pi}\int_{U'}g_{V'}(\theta)\,d\theta .
\]
Let $\Omega=\T\setminus V'$ and let $t_0$ be such that the set $U^\star=\{\theta\in\Omega:g_{V'}(\theta)>t_0\}$
has measure $\alpha$; such $t_0$ exists because $g_{V'}$ is real analytic and non-constant on each arc of $\Omega$,
so its level sets are finite and the distribution function is continuous. Then $g_{V'}\le t_0$ on
$U'\setminus U^\star$, $g_{V'}\ge t_0$ on $U^\star\setminus U'$, and $|U'\setminus U^\star|=|U^\star\setminus U'|$, so
\[
\int_{U'}g_{V'}-\int_{U^\star}g_{V'}=\int_{U'\setminus U^\star}g_{V'}-\int_{U^\star\setminus U'}g_{V'}\le0
\]
(the bathtub principle, cf.~\cite[Thm.~1.14]{LL}). The value $\int_{U^\star}g_{V'}=\int_0^{\alpha}g^*(s)\,ds$, with $g^*$
the decreasing rearrangement of $g_{V'}|_\Omega$, depends only on the distribution function of $g_{V'}$ on
$\Omega$, hence, by Lemma~\ref{lem:boole}, only on $\alpha$ and $\beta$. For the arc $V'_0=(-\beta,0)$,
\[
g_{V'_0}(\theta)=\frac1\pi\log\frac{\sin\frac{\theta+\beta}2}{\sin\frac\theta2},\qquad
g_{V'_0}'(\theta)=\frac1{2\pi}\Bigl[\cot\frac{\theta+\beta}2-\cot\frac\theta2\Bigr]<0\qquad(0<\theta<2\pi-\beta),
\]
so its superlevel sets in $\T\setminus V'_0=(0,2\pi-\beta)$ are initial intervals and $U^\star=(0,\alpha)$. Hence
$\int_{U'}g_{V'}\le\int_0^\alpha g_{V'_0}$, which after rescaling is the claim.
\end{proof}

\begin{remark}[degeneracy of the continuous problem]\label{rem:degenerate}
The continuous inequality has many maximisers. For an integer $w\ge1$ let $\pi_w(\theta)=w\theta$ on $\T$. The
conjugate function commutes with $f\mapsto f\circ\pi_w$ (the multiplier $-i\,\mathrm{sgn}(n)$ is invariant under
$n\mapsto wn$), and $\pi_w$ preserves the normalised measure, so the $w$ arcs $\pi_w^{-1}((-\beta,0))$ and the $w$
arcs $\pi_w^{-1}((0,\alpha))$ give the same value $\int_0^\alpha g_{V'_0}$. In the discrete problem these
correspond to cyclically ordered colourings with winding number $w$, which are therefore excluded only by the
junction terms $E(A,B)$ of Lemma~\ref{lem:smear}.
\end{remark}

\begin{remark}[the discrete problem needs more]\label{rem:discrete}
(i) The discrete analogue of Lemma~\ref{lem:boole} fails: for $N=7$ the values of $p_B$ on $\Z_7\setminus B$ are
$\pm2.874,\pm1.026,0$ for $B=\{0,1\}$, but $\pm2.305,\pm0.569,0$ for $B=\{0,2\}$ (to three decimals). This is why
the proof passes through the smeared sets and pays for the discretisation separately. (ii) Theorem~\ref{thm:A} is
special to the cotangent kernel. For $N=20$ and the odd kernel $\kappa'=\one_{\{1\}}-\one_{\{-1\}}$, which is also
non-increasing on $\{1,\dots,N-1\}$, the pair $A=\{0,10\}$, $B=\{9,19\}$ gives
$\sum_{x\in A,y\in B}\kappa'(x-y)=2$, while two adjacent intervals of sizes $2$ and $2$ give $1$.
\end{remark}

\subsection{The junction function}
For real $u\ge1$ let
\[
\varepsilon(u)=\frac1u-\int_{-1}^{1}\frac{1-|v|}{u+v}\,dv
=\frac1u-\bigl[(u+1)\log(u+1)-2u\log u+(u-1)\log(u-1)\bigr]
\]
(with $0\log0=0$), and $\varepsilon(-u)=-\varepsilon(u)$. The second expression holds because
$\int_{-1}^1(1-|v|)f(u+v)\,dv$ is the second difference $G(u+1)-2G(u)+G(u-1)$ of any $G$ with $G''=f$, and
$G(x)=x\log x$ for $f(x)=1/x$ (at $u=1$ by continuity). In particular $\varepsilon(1)=1-2\log2$.

\begin{lemma}[junction function]\label{lem:junction}
\begin{enumerate}
\item For $u\ge1$, $\displaystyle\varepsilon(u)=-\sum_{i\ge2}\frac{u^{-(2i-1)}}{i(2i-1)}$. Hence $\varepsilon(u)<0$,
$|\varepsilon|$ is strictly decreasing on $[1,\infty)$, and for $u\ge2$
\[
|\varepsilon(u)|\le\frac1{6u^3}+\frac{4}{45\,u^5}.
\]
\item For $m\in\{1,\dots,N-1\}$, $e(m)=\frac N\pi\sum_{j\in\Z}\varepsilon(m+jN)$, the series converging
absolutely.
\item $e(N-m)=-e(m)$. For $1\le m\le N/2$, $\frac N\pi\varepsilon(m)\le e(m)\le0$, so $|e(m)|\le\frac N\pi|\varepsilon(m)|$.
Moreover $|e(1)|\ge\frac N\pi\bigl(|\varepsilon(1)|-|\varepsilon(N-1)|\bigr)$, and for $N\ge6$
\[
|e(1)|\ge 0.38493\,\frac N\pi .
\]
\item For $N\ge 4$ put $\sprod=\sum_{m=2}^{\lfloor N/2\rfloor}(m-1)|e(m)|$ and $\fw=\sum_{m=2}^{\lfloor N/2\rfloor}m\,|e(m)|$.
Then
\begin{align*}
\sprod&\le\Bigl[\frac{\zeta(2)-\zeta(3)}6+\frac4{45}\bigl(\zeta(4)-\zeta(5)\bigr)\Bigr]\frac N\pi\le0.07785\,\frac N\pi,\\
\fw&\le\Bigl[\frac{\zeta(2)-1}6+\frac4{45}\bigl(\zeta(4)-1\bigr)\Bigr]\frac N\pi\le0.11481\,\frac N\pi .
\end{align*}
\end{enumerate}
\end{lemma}

\begin{proof}
(1) For $u>1$ and $|v|\le1$, $\frac1{u+v}=\sum_{n\ge0}(-v)^nu^{-n-1}$ uniformly in $v$. Odd powers integrate to zero
against $1-|v|$, and $\int_{-1}^1(1-|v|)v^{2i}dv=\frac2{(2i+1)(2i+2)}$, so
$\int_{-1}^1\frac{1-|v|}{u+v}dv=\sum_{i\ge0}\frac{u^{-2i-1}}{(i+1)(2i+1)}$. The term $i=0$ is $1/u$; shifting the index
gives the series for $\varepsilon(u)$. At $u=1$ it follows by monotone convergence as $u\downarrow1$ (all terms are
positive and both sides are continuous). Each term of $|\varepsilon(u)|$ is strictly decreasing in $u$. For $u\ge2$
the term $i=2$ is $1/(6u^3)$, and for $i\ge3$ we have $i(2i-1)\ge15$ and $u^{-(2i-1)}\le u^{-5}4^{-(i-3)}$, so the
remaining terms add up to at most $\frac{u^{-5}}{15}\sum_{j\ge0}4^{-j}=\frac{4}{45}u^{-5}$. (The exact value of
$\sum_{i\ge3}4^{3-i}/(i(2i-1))$ is $0.07727\ldots<0.0773<4/45$.)

(2) The partial fraction expansion $\pi\cot\pi s=\lim_{J\to\infty}\sum_{j=-J}^J\frac1{s+j}$, locally uniform on
$\R\setminus\Z$, gives $\kappa(t)=\frac N\pi\lim_J\sum_{|j|\le J}\frac1{t+jN}$. On $[m-1,m+1]$ a term can have a pole
only at an endpoint (the term $j=0$ at $t=0$ if $m=1$, the term $j=-1$ at $t=N$ if $m=N-1$); such terms are
integrable against the tent weight, which vanishes at the endpoints, and the sum of the remaining terms converges
uniformly on $[m-1,m+1]$. Integrating term by term gives $\frac\pi Ne(m)=\lim_J\sum_{|j|\le J}\varepsilon(m+jN)$, because for every real $u$ with $|u|\ge1$ one has
$\frac1u-\int_{-1}^1\frac{1-|v|}{u+v}dv=\varepsilon(u)$ (for $u\le-1$ substitute $v\mapsto-v$ and use that $\varepsilon$ is odd).
The series converges absolutely by (1), since $|\varepsilon(u)|\le\frac1{6|u|^3}+\frac4{45|u|^5}$ for $|u|\ge2$.

(3) Both $\kappa$ and $\kb$ are odd under $m\mapsto N-m$ (for $\kb$ substitute $v\mapsto-v$). Let $1\le m\le N/2$. Grouping the
terms of (2) as $j=0,-1,1,-2,2,\dots$ gives
\[
\frac\pi Ne(m)=-|\varepsilon(m)|+|\varepsilon(N-m)|-|\varepsilon(N+m)|+|\varepsilon(2N-m)|-\cdots,
\]
an alternating series whose arguments $m\le N-m\le N+m\le2N-m\le\cdots$ are non-decreasing, so its terms are
non-increasing in absolute value and tend to zero. Hence the sum lies between the first two partial sums:
$-|\varepsilon(m)|\le\frac\pi Ne(m)\le-|\varepsilon(m)|+|\varepsilon(N-m)|\le0$. For $m=1$ the right inequality gives the
lower bound for $|e(1)|$. Numerically $|\varepsilon(1)|=2\log2-1=0.3862943\ldots$ and, for $N\ge6$,
$|\varepsilon(N-1)|\le|\varepsilon(5)|=|\tfrac15-(6\log6-10\log5+4\log4)|=0.0013551\ldots$, so
$|\varepsilon(1)|-|\varepsilon(N-1)|\ge0.3849392\ldots\ge0.38493$.

(4) By (3) and (1), $\sprod\le\frac N\pi\sum_{m\ge2}(m-1)\bigl(\frac1{6m^3}+\frac4{45m^5}\bigr)$ and
$\fw\le\frac N\pi\sum_{m\ge2}m\bigl(\frac1{6m^3}+\frac4{45m^5}\bigr)$; the sums are the stated combinations of
$\zeta(2)=\pi^2/6$, $\zeta(3)$, $\zeta(4)=\pi^4/90$, $\zeta(5)$, whose values are $0.0778480\ldots$ and
$0.1148066\ldots$.
\end{proof}

The numerical constants are rounded in the safe direction; they are re-derived in interval arithmetic by the
script \texttt{check\_constants.py} (Section~\ref{sec:numerics}). The proof only needs
$|e(1)|>2\,\sprod+\fw$, and the bounds give $0.38493>2\cdot0.07785+0.11481=0.27051$ (in units of $N/\pi$); see
\eqref{eq:margin}.

\subsection{The junction inequality}

\begin{lemma}[junction inequality]\label{lem:junctionineq}
Let $a,b,c\ge1$ and let $(A,B)$ be cyclically ordered with winding number $w\ge2$. Put $A_0=[0,a)$, $B_0=[-b,0)$.
Then $N\ge6$ and
\begin{equation}\label{eq:margin}
E(A,B)-E(A_0,B_0)\le-(w-1)|e(1)|+w\,\sprod+\fw\le\frac N\pi\,(0.49974-0.30708\,w)<0 .
\end{equation}
\end{lemma}

\begin{proof}
Since there are $3w\ge6$ nonempty blocks, $N\ge6$. Sort the pairs $(x,y)\in A\times B$ by $m=x-y\in\{1,\dots,N-1\}$.

\emph{$m=1$:} then $y\in B$ and $y+1=x\in A$, a colour change $B\to A$; in a cyclically ordered colouring these are
exactly the $w$ boundaries between $B_i$ and $A_i$. They contribute $-w|e(1)|$.

\emph{$m=N-1$:} then $x\in A$ and $x+1=y\in B$, a change $A\to B$, which does not occur.

\emph{$2\le m\le N/2$:} these pairs contribute $\le0$ by Lemma~\ref{lem:junction}(3).

\emph{$N/2<m\le N-2$:} put $m'=N-m=y-x\bmod N\in[2,N/2)$; the pair contributes $e(m)=|e(m')|$. Let $x\in A_i$ and let
$r$ be the number of points of $A_i$ after $x$. Going forward from $x$ one meets these $r$ points, then the nonempty
block $C_{i+1}$, and only then points of $B$; so $m'\ge r+2$. For fixed $x$ and $m'$ there is at most one $y$. Summing
over the points of $A_i$, for which $r=0,1,\dots,|A_i|-1$, the contribution of $A_i$ is at most
\[
\sum_{r\ge0}\ \sum_{m'=r+2}^{\lfloor N/2\rfloor}|e(m')|=\sum_{m'=2}^{\lfloor N/2\rfloor}(m'-1)|e(m')|=\sprod,
\]
and the total over the $w$ blocks is at most $w\,\sprod$. Hence $E(A,B)\le-w|e(1)|+w\,\sprod$.

For the pair $(A_0,B_0)$ all differences $m=x-y$ lie in $[1,a+b-1]\subseteq[1,N-2]$, the value $m=1$ occurs once,
and each value $m$ occurs at most $m$ times, because $x\in[0,a)$ and $y=x-m\in[-b,0)$ force $0\le x\le m-1$ and $y$
is determined by $x$. Pairs with $m>N/2$ contribute $\ge0$ and pairs with $2\le m\le N/2$ contribute
$\ge-m|e(m)|$ in total for each $m$. Hence $E(A_0,B_0)\ge-|e(1)|-\fw$. Subtracting gives the first inequality in \eqref{eq:margin};
inserting $|e(1)|\ge0.38493\,N/\pi$, $\sprod\le0.07785\,N/\pi$ and $\fw\le0.11481\,N/\pi$ from
Lemma~\ref{lem:junction} gives $\frac N\pi\bigl(-(w-1)0.38493+0.07785\,w+0.11481\bigr)=\frac N\pi(0.49974-0.30708\,w)$,
which is at most $-0.11442\,N/\pi<0$ for $w\ge2$.
\end{proof}

\subsection{Conclusion of the proof}
\begin{proof}[Proof of Theorem~\ref{thm:A}]
The trivial cases were treated at the beginning of this section; let $a,b,c\ge1$. There are finitely many pairs
$(A,B)$, so a maximiser exists. By Corollary~\ref{cor:monotone} it is cyclically ordered with some winding number
$w\ge1$. Suppose $w\ge2$. By Lemma~\ref{lem:smear}, Proposition~\ref{prop:cont} (applied to $U=\widetilde A$,
$V=\widetilde B$, which are disjoint finite unions of half-open intervals of measures $a$ and $b$, and
$\widetilde{A_0}=[0,a)$, $\widetilde{B_0}=[-b,0)$) and Lemma~\ref{lem:junctionineq},
\[
\ip AB=\Qc(\widetilde A,\widetilde B)+E(A,B)\le\Qc(\widetilde{A_0},\widetilde{B_0})+E(A,B)
<\Qc(\widetilde{A_0},\widetilde{B_0})+E(A_0,B_0)=\ip{A_0}{B_0},
\]
contradicting maximality. So $w=1$, the maximiser is a rotation of $(A_0,B_0)$ by Corollary~\ref{cor:monotone},
and the maximum is $\Phi_N(a,b)$ because $\ip{\cdot}{\cdot}$ is invariant under rotations. Every pair attaining
$\Phi_N(a,b)$ is a maximiser, hence a rotation of $(A_0,B_0)$.
\end{proof}

\section{Numerical cross-checks}\label{sec:numerics}

The proofs above are complete without computer assistance. The following checks were run to guard against
errors in formulas and signs; they are \emph{numerical} (floating point, except where interval arithmetic is
stated) and are not part of the proofs. The scripts and their outputs are in the folders \texttt{scripts/} and
\texttt{logs/} of the accompanying material; \texttt{scripts/README.md} lists the commands.

\begin{table}[h]
\small
\begin{tabular}{@{}>{\raggedright\arraybackslash}p{0.27\textwidth}>{\raggedright\arraybackslash}p{0.21\textwidth}>{\raggedright\arraybackslash}p{0.46\textwidth}@{}}
\toprule
Statement checked & Script & Range and outcome\\
\midrule
Theorem~\ref{thm:B} for $M=1$, by exhaustion & \texttt{verify\_writeup.py} & all $4^d$ configurations, $d=2,\dots,9$:
$\max F=\FD$ (difference $\le 3\cdot10^{-14}$), exactly $4d$ maximisers, all intervals\\
Lemma~\ref{lem:cot} & \texttt{verify\_writeup.py}, \texttt{cot\_rearr.py} & identity error $\le 8\cdot10^{-14}$, $d=2,\dots,9$\\
Theorem~\ref{thm:A} by exhaustion & \texttt{cot\_rearr.py} & $(N,a,b)=(12,3,3)$, $(16,4,4)$, $(20,5,5)$, $(24,6,6)$,
$(16,3,6)$, $(20,4,7)$: maximum $=\Phi_N(a,b)$, attained by adjacent intervals\\
Lemmas~\ref{lem:cyclic}, \ref{lem:swap}, \ref{lem:smear} & \texttt{verify\_writeup.py} & random colourings,
$N\in\{8,12,20,33,40\}$: errors $\le4\cdot10^{-13}$\\
Corollary~\ref{cor:monotone} (exchanges) & \texttt{swap\_test.py} & all colourings for $(N,a,b)=(9,3,3)$, $(10,3,3)$,
$(12,3,4)$, $(12,4,4)$: $290{,}710$ out-of-order exchanges, none decreases $\ip AB$\\
Lemma~\ref{lem:junction}(2) & \texttt{verify\_writeup.py} & $N\le40$: agreement to $2\cdot10^{-7}$ (quadrature accuracy)\\
Lemma~\ref{lem:junction}(3),(4) & \texttt{verify\_junction.py} & $N\in\{6,7,8,12,16,24,40,64,101\}$: signs and bounds hold;
$0.38536\le|e(1)|/(N/\pi)\le0.38630$, $\sprod/(N/\pi)\le0.0733$, $\fw/(N/\pi)\le0.1097$\\
constants in Lemmas~\ref{lem:junction} and~\ref{lem:junctionineq} & \texttt{check\_constants.py} & interval arithmetic
(60 digits): all stated roundings are valid\\
whole chain on exchange-flow endpoints & \texttt{verify\_chain.py} & random starts, $N\le80$: exchange identity to
$3\cdot10^{-13}$; $\Qc\le\Qc^{\mathrm{block}}$, $E\le E^{\mathrm{block}}$ and $\ip AB\le\Phi_N$ at every endpoint\\
Section~\ref{sec:model}, Appendix~\ref{app:reduction} & \texttt{check\_reduction.py} & checks (R1)--(R9) below\\
Sherali--Adams relaxation & \texttt{classical\_sa.py} & $d=3,\dots,12$ (Section~\ref{sec:discussion})\\
\bottomrule
\end{tabular}
\end{table}

The script \texttt{check\_reduction.py} computes CGLMP values from outcome probabilities with the expression
\eqref{eq:cglmp} and checks: (R1) Lemma~\ref{lem:chain} on random behaviours; (R2) formula \eqref{eq:Strace} on
random projective strategies on $\Phi_D$, $D\le6$; (R3) Theorem~\ref{thm:reduction}(1) for random families $V$ with
$M\le3$ (non-commuting for $M\ge2$); (R4) Theorem~\ref{thm:reduction}(2), by carrying out the twirl and the Stone--von Neumann
decomposition of Appendix~\ref{app:reduction} on random strategies; (R5) Proposition~\ref{prop:links} on random
strategies (unequal links, non-commuting reduced family) and on equal-link strategies (commuting reduced family);
(R6) exhaustion over all clock-phase strategies for $d\le7$ (maximum $\IME(d)$, attained exactly by the $4d$
one-step ones); (R7) the value $\IME(d)$ of the CGLMP measurements and the identity $R^{\mathrm{DKZ}}=GR^{(0)}G^*$ of
Section~\ref{subsec:dkz}, $d\le8$; (R8) the projection form \eqref{eq:projform} and (R9) the Ising form of
Section~\ref{sec:discussion}, for random non-commuting families. All identities hold to within $2\cdot10^{-14}$.

\section{Discussion}\label{sec:discussion}

\subsection{An Ising reformulation}
A unitary $E$ with $E^4=1$ can be written as $E=e^{i\pi/4}(X-iX')/\sqrt2$ with commuting involutions $X,X'$: on
the eigenspaces of $E$ for the eigenvalues $1$, $i$, $-1$, $-i$ take $(X,X')=(1,1)$, $(1,-1)$, $(-1,-1)$, $(-1,1)$,
respectively. Put $E_k=V_k^*$,
$X_k=X(E_k)$ and $X_{k+d}=X'(E_k)$ for $0\le k\le d-1$, and extend antiperiodically by $X_{n+2d}=-X_n$. If one
expands $\tau(E_j^*E_k)$ and groups the terms by their distance on the ring $\Z_{2d}$, one obtains
\[
F(V)=\frac12\sum_{n\in\Z_{2d}}\sum_{u=1}^{d-1}\sec\frac{\pi u}{2d}\;\tau(X_nX_{n+u}),
\]
and conversely every family of involutions with $[X_n,X_{n+d}]=0$ arises. In the commutative sector this is the
energy of a classical long-range ferromagnetic Ising ring of $2d$ spins with antiperiodic boundary condition;
the DKZ configuration is a single domain wall. Theorem~\ref{thm:B} thus says that for these couplings the
single domain wall is the unique ground state (up to translation and global sign). The constraint
$[X_n,X_{n+d}]=0$ matters: the anticommuting choice $X_n=\cos(\frac{\pi n}{2d})\sigma_1+\sin(\frac{\pi n}{2d})\sigma_2$
(Pauli matrices, $\tau=\frac12\Tr$) gives $\tau(X_nX_{n+u})=\cos\frac{\pi u}{2d}$ and hence the value $d(d-1)$, which
exceeds $\FD$ because $(d-m)\sec\frac{\pi m}{2d}<d$ for $1\le m\le d-1$ (use $\cos x>1-\frac{2x}\pi$ on $(0,\frac\pi2)$).

\subsection{Why local certificates are not enough}
The degeneracy in Remark~\ref{rem:degenerate} is visible in the clock model itself. Let $w\ge5$ with
$w\equiv1\pmod4$ divide $d$, and let $A=\bigcup_{j=0}^{w-1}\bigl(\tfrac{4d}{w}j+[0,\tfrac dw)\bigr)\subseteq\Z_{4d}$. Since
$d\equiv\frac dw\pmod{\frac{4d}w}$, the translates $A,A+d,A+2d,A+3d$ are the $w$-fold unions of the four intervals of
length $\frac dw$ in each period $\frac{4d}{w}$, so they partition $\Z_{4d}$; hence $A$ is a transversal, $A=S_a$ for some
$a\in\Z_4^d$, and $(\widetilde A,\widetilde{A-d})$ is the $w$-fold cover of $([0,d),[-d,0))$. By
Remark~\ref{rem:degenerate}, Lemma~\ref{lem:smear} and Lemma~\ref{lem:cot},
\[
\FD-F(a)=\tfrac12\bigl[E(A_0,B_0)-E(A,A-d)\bigr]\qquad(A_0=[0,d),\ B_0=[-d,0)),
\]
a difference of junction terms. The pair $(A,A-d)$ is cyclically ordered with winding number $w$, so by
Lemma~\ref{lem:junctionineq} this difference is at least $\frac{2d}\pi(0.30708\,w-0.49974)$. It is also at most
$\frac{2d}\pi\,0.50111\,w$: all differences occurring in $(A_0,B_0)$ lie in $[1,2d-1]$, so $E(A_0,B_0)\le0$ by
Lemma~\ref{lem:junction}(3); and in $(A,A-d)$ a point $y$ of a block $B_i$ followed by $s$ further points of $B_i$
has its first $A$-point at distance $s+1$, so, counting as in the proof of Lemma~\ref{lem:junctionineq},
$E(A,A-d)\ge-w\sum_{m=1}^{2d}m|e(m)|=-w(|e(1)|+\fw)\ge-w(0.38630+0.11481)\frac{4d}\pi$. So the gap is of order $d$
for fixed $w$, while $\FD\ge d(d-1)/2$: such ``$w$-fold'' configurations are separated from the DKZ
configuration only at relative order $1/d$, and one should not expect certificates built from the marginals of a
bounded number of sites to resolve this for large $d$. As an illustration (numerical, linear programming with the
HiGHS solver, script
\texttt{classical\_sa.py}): the relaxation of the classical clock model in which one optimises over consistent
families of three-site marginals of the differences $a_k-a_j$ is exact for $d=3,\dots,8$, but exceeds $\FD$ by
$0.535$, $2.52$, $5.23$ and $8.53$ for $d=9,10,11,12$. The argument of Section~\ref{sec:A} is global: it uses the
cyclic order of the whole configuration, and the junction estimate compares different windings directly.

\subsection{The non-commutative case}
The first steps of the argument have operator versions. Let $V_0,\dots,V_{d-1}$ be arbitrary (not necessarily
commuting) unitaries of order dividing four in a tracial von Neumann algebra $(\mathcal N,\tau)$, $E_k=V_k^*$, and for
$x\in\Z_{4d}$ let $Q_x$ be the spectral projection of $E_k$ for the eigenvalue $i^a$, where $x=k-da$ ($0\le k\le d-1$,
$a\in\Z_4$). Put $A_x=Q_x$, $B_x=Q_{x+d}$ and $N=4d$. Then $A_xB_x=0$, $\sum_xA_x=\sum_xB_x=d\cdot\one$, and the
computation in the proof of Lemma~\ref{lem:cot} gives
\begin{equation}\label{eq:projform}
F(V)=\frac12\sum_{x,y\in\Z_{4d}}\kappa(x-y)\,\tau(A_xB_y)-\frac d2 .
\end{equation}
This suggests the following operator version of Theorem~\ref{thm:A}.

\begin{question}\label{q:Lq}
Let $A_x,B_x$ ($x\in\Z_N$) be orthogonal projections in a tracial von Neumann algebra with $A_xB_x=0$ for all
$x$, $\sum_xA_x=a\cdot\one$ and $\sum_xB_x=b\cdot\one$. Is
$\sum_{x,y}\kappa(x-y)\,\tau(A_xB_y)\le\Phi_N(a,b)$?
\end{question}

A positive answer for $N=4d$ and $a=b=d$ (even only for the families of the form above) would prove the optimality
of the DKZ measurements on maximally entangled states for all $d$. For commuting families the answer is positive by
Theorem~\ref{thm:A}, applied to each joint eigenvector. For matrices the same holds if the $A_x$ commute among
themselves and the $B_x$ commute among themselves: if $(e_i)$ and $(f_j)$ are joint eigenbases, with
$A_xe_i=[x\in S_i]e_i$ and $B_xf_j=[x\in T_j]f_j$, then $|S_i|=a$, $|T_j|=b$, the matrix $|\langle e_i,f_j\rangle|^2$ is
doubly stochastic, $\langle e_i,f_j\rangle=0$ unless $S_i\cap T_j=\emptyset$ (because $A_xB_x=0$), and the left side
equals $\frac1M\sum_{i,j}|\langle e_i,f_j\rangle|^2\ip{S_i}{T_j}\le\Phi_N(a,b)$. The operator constraints cannot be
weakened to $\sum_x\tau(A_x)=a$, $\sum_x\tau(B_x)=b$: a mixture that, with probability $\frac12$, uses the sets
$([0,2),[-2,0))$ and otherwise empty sets has $\sum_x\tau(A_x)=\sum_x\tau(B_x)=1$ and value
$\frac12\Phi_{20}(2,2)=7.216\ldots>\Phi_{20}(1,1)=\cot(\pi/20)=6.313\ldots$ for $N=20$.

The cyclic symmetry of Lemma~\ref{lem:cyclic} holds verbatim for such families, with the projections
$C_x=\one-A_x-B_x$ and $\ip AB_\tau=\sum_{x,y}\kappa(x-y)\tau(A_xB_y)$. The exchange identity of Lemma~\ref{lem:swap}
survives in a weak form: if $X$ is a projection with $X\le A_x$ and $X\le B_{x+1}$, then replacing
$(A_x,B_x,A_{x+1},B_{x+1})$ by $(A_x-X,B_x+X,A_{x+1}+X,B_{x+1}-X)$ preserves all constraints, leaves every $C_t$
unchanged, and changes the pairing by $\sum_{t\ne x,x+1}\delta(x-t)\,\tau(XC_t)\ge0$ (use $\ip AB_\tau=\ip BC_\tau$).
In the non-commutative case such subprojections need not exist, and we do not know operator analogues of
Proposition~\ref{prop:cont} and of the junction inequality.

\subsection{A non-commuting class: two-window families}
There is one class of genuinely non-commuting families where Theorem~\ref{thm:B} extends directly. In it, every
unitary switches between the values of two optimal configurations.

\begin{theorem}[two-window families]\label{thm:D}
Let $a,a'\in\Z_4^d$ be one-step configurations. Let $P_0,\dots,P_{d-1}$ be arbitrary orthogonal projections in a von
Neumann algebra with faithful normal tracial state $\tau$ (for instance $M_M(\C)$ with $\tr$), and put
\[
V_k=i^{-a_k}(\one-P_k)+i^{-a'_k}P_k\qquad(0\le k\le d-1).
\]
Then $F(V)\le\FD$. Equality holds if and only if the projections $P_k$ with $a_k\ne a'_k$ are totally ordered (so that
all $V_k$ commute) and every joint eigen-configuration is one-step.
\end{theorem}

\begin{proof}
Fix $j<k$, put $m=k-j$, $\psi=\psi_m$, and let $c_t=\Re[h_mi^t]$, so that $(c_0,c_1,c_2,c_3)=(\sec\psi,\csc\psi,-\sec\psi,-\csc\psi)$.
With $p_k=\tau(P_k)$ and $q_{jk}=\tau(P_jP_k)=\tau(P_jP_kP_j)\in[0,\min(p_j,p_k)]$,
\[
\tau(V_jV_k^*)=i^{a_k-a_j}(1-p_j-p_k+q_{jk})+i^{a'_k-a_j}(p_k-q_{jk})+i^{a_k-a'_j}(p_j-q_{jk})+i^{a'_k-a'_j}q_{jk},
\]
so $\Re[h_m\tau(V_jV_k^*)]=\ell_{jk}(p_j,p_k)+w_{jk}q_{jk}$, with $\ell_{jk}$ affine and
$w_{jk}=c_{a_k-a_j}+c_{a'_k-a'_j}-c_{a'_k-a_j}-c_{a_k-a'_j}$. If $a_j=a'_j$ or $a_k=a'_k$ then $w_{jk}=0$.

Otherwise, let $\varepsilon=a_k-a_j$ and $\varepsilon'=a'_k-a'_j$, both in $\{0,1\}$ since the configurations are
one-step and $j<k$, and let $\delta=a'_j-a_j\in\{1,2,3\}$. The condition $a_k\ne a'_k$ reads
$\delta+\varepsilon'-\varepsilon\ne0$. Evaluating the four cases:
\begin{itemize}
\item $(\varepsilon,\varepsilon')=(0,0)$: $w_{jk}=2c_0-c_\delta-c_{-\delta}$, which is $2\sec\psi$ for $\delta$ odd and $4\sec\psi$ for
$\delta=2$.
\item $(1,1)$: $w_{jk}=2c_1-c_{1+\delta}-c_{1-\delta}$, which is $2\csc\psi$ for $\delta$ odd and $4\csc\psi$ for $\delta=2$.
\item $(0,1)$ with $\delta\in\{1,2\}$, and $(1,0)$ with $\delta\in\{2,3\}$: $w_{jk}=c_0+c_1-c_{\delta+\varepsilon'}-c_{\varepsilon-\delta}=2(\sec\psi+\csc\psi)$.
\end{itemize}
Hence $w_{jk}\ge2\min(\sec\psi_m,\csc\psi_m)>0$ for every such pair.

Consequently $F(V)=L(p)+\sum w_{jk}q_{jk}\le L(p)+\sum w_{jk}\min(p_j,p_k)$, where $L$ collects the affine terms and the
sums run over the pairs with $w_{jk}\ne0$. The right-hand side is $F(V')$ for the family $V'_k=i^{-a_k}(\one-P'_k)+i^{-a'_k}P'_k$
with the nested projections $P'_k=\mathbf 1_{[0,p_k)}$ in $L^\infty[0,1]$, which has the same traces and
$\tau(P'_jP'_k)=\min(p_j,p_k)$. The family $V'$ commutes, so $F(V')\le\FD$ by Theorem~\ref{thm:B}.

For equality, $q_{jk}=\min(p_j,p_k)$ is needed for every pair with $w_{jk}>0$. If $p_j\le p_k$, this gives
$\tau(P_j(\one-P_k)P_j)=0$, hence $P_j\le P_k$ because $\tau$ is faithful. So the $P_k$ with $a_k\ne a'_k$ form a
chain, all $V_k$ commute, and the equality statement of Theorem~\ref{thm:B} applies. The converse is clear.
\end{proof}

For $M\ge2$ the families of Theorem~\ref{thm:D} are in general far from commuting: the projections $P_k$ are arbitrary. The
proof shows that on this class $F$ is exactly affine in the Gram data $\tau(P_jP_k)$ with positive coefficients. This
turns the local (second-order) optimality of DKZ, whose Hessian involves the same coefficients~\cite{repo}, into a
global statement. For three or more values per unitary, the corresponding coefficients are no longer all positive,
so the argument does not extend directly. The script \texttt{check\_two\_window.py} confirms numerically that the smallest
coefficient equals $2\sec(\pi/2d)$ for $2\le d\le30$, and that $F(V)\le\FD$ for random two-window families.

\subsection{Summary of the status of Problem 27B (maximally entangled part)}
The inequality $F(V)\le\FD$ (equivalently $I_d\le\IME(d)$ for projective measurements on maximally entangled
states) is proved for commuting families and for the non-commuting two-window families of Theorem~\ref{thm:D}, for
all $d$ (this paper), and for all families for $d\le20$ by exact certificates~\cite{repo}; for $d=2$ it is
Tsirelson's bound. It is open for general non-commuting families when $d\ge21$.
Within the commutative sector, Corollary~\ref{cor:C} also gives the uniqueness (``necessarily of DKZ form'') asked
for in~\cite{OQP}.

\appendix
\section{Proof of Theorem~\ref{thm:reduction} and Proposition~\ref{prop:links}}\label{app:reduction}

Write $\tau$ for the normalised trace on the relevant matrix algebra.

\emph{Invariances.} For a quadruple $R=(R_1,R_2,R_3,R_4)$ of unitaries with $R_i^d=1$ let
$\rho(R)=(R_2,R_3,R_4,wR_1)$, again a quadruple of this kind. Directly from the definitions,
$L^{(0)}_n(\rho R)=R_2^nR_3^{-n}=L^{(1)}_n(R)$, $L^{(1)}_n(\rho R)=L^{(2)}_n(R)$,
$L^{(2)}_n(\rho R)=R_4^n(wR_1)^{-n}=L^{(3)}_n(R)$ and $L^{(3)}_n(\rho R)=w^{-n}(wR_1)^nR_2^{-n}=L^{(0)}_n(R)$. So
$L^{(i)}_n(\rho R)=L^{(i+1)}_n(R)$ with the upper index read modulo $4$; in particular $\Lambda_n(\rho R)=\Lambda_n(R)$ and
$S(\rho R)=S(R)$ by \eqref{eq:Strace}. Moreover $\rho^4(R)=wR$ and $\rho^{4d}(R)=R$. For a direct sum $R\oplus R'$ on
$\C^D\oplus\C^{D'}$, $S$ is the average of $S(R)$ and $S(R')$ with weights $D$ and $D'$.

\emph{Covariant strategies: proof of (1).} Let $R=R^V$ as in \eqref{eq:covariant}. From $W^4=\sum_k|k\rangle\langle k|\otimes w^k=Z_0\otimes\one$
and $Z_0XZ_0^*=wX$ we get $WR_4W^*=W^4R_1W^{-4}=wR_1$. Hence $L^{(1)}_n=WL^{(0)}_nW^*$, $L^{(2)}_n=W^2L^{(0)}_nW^{-2}$, and
$W^3L^{(0)}_nW^{-3}=W^3R_1^nW^{-3}\,W^4R_1^{-n}W^{-4}=w^{-n}R_4^nR_1^{-n}=L^{(3)}_n$. So all four links have the trace
$\tau(L^{(0)}_n)$. Since $R_1^nWR_1^{-n}=\sum_k|k\rangle\langle k|\otimes W_{k-n}$ with $W_k=z^kV_k$ and indices modulo $d$,
\[
L^{(0)}_n=\sum_{k}|k\rangle\langle k|\otimes W_{k-n}W_k^*,\qquad
\tau(L^{(0)}_n)=\frac{z^{-n}}d\Bigl[\sum_{k=n}^{d-1}\tr(V_{k-n}V_k^*)+i\sum_{k=0}^{n-1}\tr(V_{k-n+d}V_k^*)\Bigr],
\]
using $z^d=i$ for $k<n$. Since $w^{-n}=e^{-4i\psi_n}$ and $z^{-n}=e^{-i\psi_n}$, we have
$c_n=-1/(1-e^{-4i\psi_n})=\frac i2\,e^{2i\psi_n}/\sin(2\psi_n)$ and
$4c_nz^{-n}=i\,e^{i\psi_n}/(\sin\psi_n\cos\psi_n)=i\csc\psi_n-\sec\psi_n=-h_n$. Therefore
\[
\sum_{n=1}^{d-1}c_n\Lambda_n=4\sum_nc_n\tau(L^{(0)}_n)
=-\frac1d\sum_{n}h_n\Bigl[\sum_{k\ge n}\tr(V_{k-n}V_k^*)+i\sum_{k<n}\tr(V_{k-n+d}V_k^*)\Bigr].
\]
The first inner sum runs over the pairs $j=k-n<k$ with $k-j=n$ and gives $\sum_{j<k}h_{k-j}\tr(V_jV_k^*)$. In the second,
put $p=k<q=k-n+d$, so $q-p=d-n$; since $\psi_{d-n}=\frac\pi2-\psi_n$ we have $h_{d-m}=\csc\psi_m-i\sec\psi_m$, i.e.\
$i\,h_{d-m}=\overline{h_m}$, and the second part equals $\sum_{p<q}\overline{h_{q-p}}\,\tr(V_qV_p^*)
=\sum_{p<q}\overline{h_{q-p}\tr(V_pV_q^*)}$. Hence $\sum_nc_n\Lambda_n=-\frac2dF(V)$ and
$S(R^V)=2(d-1)-\frac2dF(V)$; Lemma~\ref{lem:chain} gives $I_d=4F(V)/(d(d-1))$.

\emph{Twirl: proof of (2).} Let $R$ act on $\C^D$ and let $\widetilde R_i=\bigoplus_{r\in\Z_{4d}}(\rho^rR)_i$ on
$\C^{4d}\otimes\C^D$. Then $S(\widetilde R)=S(R)$, so $I_d(\widetilde R)=I_d(R)$, and $R$ is the direct summand $r=0$.
Let $W$ be the shift $W(e_{r+1}\otimes v)=e_r\otimes v$. Since $(\rho^{r+1}R)_i=(\rho^rR)_{i+1}$ for $i\le3$ and
$(\rho^{r+1}R)_4=w(\rho^rR)_1$, we get $W\widetilde R_iW^*=\widetilde R_{i+1}$ ($i\le3$), $W\widetilde R_4W^*=w\widetilde R_1$ and
$W^{4d}=1$. The unitary $Z=W^4$ satisfies $Z^d=1$ and $Z\widetilde R_1Z^*=w\widetilde R_1$. Let $H_k=\ker(Z-w^k)$,
$k\in\Z_d$; then $\widetilde R_1H_k=H_{k+1}$, all $H_k$ have the same dimension $M$, and $dM=4dD$. Choose an orthonormal
basis $(f_\mu)$ of $H_0$ and put $e_{k,\mu}=\widetilde R_1^kf_\mu$, $0\le k\le d-1$. This is an orthonormal basis in which
$\widetilde R_1=X\otimes\one_M$ (note $\widetilde R_1^d=1$) and $Z=Z_0\otimes\one_M$. Since $W$ commutes with $Z$, it
preserves every $H_k$, so $W=\sum_k|k\rangle\langle k|\otimes W_k$ with $W_k^4=w^k$; with $V_k=z^{-k}W_k$ we have $V_k^4=1$,
and $\widetilde R$ is exactly $R^V$ in this basis. Another choice of $(f_\mu)$ conjugates all $V_k$ by the same unitary.
Part (1) now gives $I_d(R)=4F(V)/(d(d-1))$. The final equivalence in the theorem follows from (1), (2) and
$\IME(d)=4\FD/(d(d-1))$.

\emph{Proof of Proposition~\ref{prop:links}.} By the computation above, $L^{(i)}_n(\rho^rR)=L^{(i+r)}_n(R)$ (upper index modulo
$4$), so $L^{(i)}_n(\widetilde R)=\bigoplus_rL^{(i+r)}_n(R)$. Hence $\widetilde R$ has equal links if and only if $R$ has.
In the basis $(e_{k,\mu})$, $\widetilde R=R^V$ and, as shown in the proof of (1), $L^{(i)}_n=W^iL^{(0)}_nW^{-i}$ for all $i$;
so $R^V$ has equal links if and only if $W$ commutes with every $L^{(0)}_n$. Comparing the diagonal blocks,
$WL^{(0)}_nW^*=\sum_k|k\rangle\langle k|\otimes W_kW_{k-n}W_k^*W_k^*$, which equals $L^{(0)}_n$ if and only if
$W_kW_{k-n}=W_{k-n}W_k$ for every $k$. For $n=1,\dots,d-1$ this says that all $W_k$, equivalently all $V_k$, commute
pairwise.

\section*{Data and code}
The check scripts, their outputs, and a description of the commands are provided with this manuscript
(folders \texttt{scripts/} and \texttt{logs/}). The exact certificates for $d\le20$ referred to above are available
at \url{https://github.com/anshM123/CGLMP}.

\begin{thebibliography}{99}

\bibitem{ADGL}
A.~Ac\'{\i}n, T.~Durt, N.~Gisin and J.~I. Latorre, \emph{Quantum nonlocality in two three-level systems},
Phys. Rev. A \textbf{65} (2002), 052325.

\bibitem{BK}
R.~Ba\~nuelos and M.~Kwa\'snicki, \emph{On the $\ell^p$-norm of the discrete Hilbert transform}, Duke Math. J.
\textbf{168} (2019), no.~3; arXiv:1709.07427.

\bibitem{Boole}
G.~Boole, \emph{On the comparison of transcendents, with certain applications to the theory of definite
integrals}, Philos. Trans. R. Soc. Lond. \textbf{147} (1857), 745--803.

\bibitem{Tsirelson}
B.~S. Cirel'son, \emph{Quantum generalizations of Bell's inequality}, Lett. Math. Phys. \textbf{4} (1980), no.~2,
93--100.

\bibitem{CGLMP}
D.~Collins, N.~Gisin, N.~Linden, S.~Massar and S.~Popescu, \emph{Bell inequalities for arbitrarily high-dimensional
systems}, Phys. Rev. Lett. \textbf{88} (2002), 040404; arXiv:quant-ph/0106024.

\bibitem{CLM}
L.~Colzani, E.~Laeng and L.~Monz\'on, \emph{Variations on a theme of Boole and Stein--Weiss}, J. Math. Anal. Appl.
\textbf{363} (2010), no.~1, 225--229.

\bibitem{DKZ}
T.~Durt, D.~Kaszlikowski and M.~\.Zukowski, \emph{Violations of local realism with quantum systems described by
$N$-dimensional Hilbert spaces up to $N=16$}, Phys. Rev. A \textbf{64} (2001), 024101.

\bibitem{GL}
C.~Giannitsi and M.~T. Lacey, \emph{Lower estimate on square function of an indicator set}, arXiv:2307.05692 (2023).

\bibitem{IR}
M.~Ioannou and D.~Rosset, \emph{Noncommutative polynomial optimization under symmetry}, arXiv:2112.10803 (2021).

\bibitem{OQP}
IQOQI Vienna, \emph{Open Quantum Problems}, Problem~27: \emph{The power of CGLMP inequalities},
\url{https://oqp.iqoqi.oeaw.ac.at/the-power-of-cglmp-inequalities}; archived copy of 29 October 2023:
\url{https://web.archive.org/web/20231029013544/https://oqp.iqoqi.oeaw.ac.at/the-power-of-cglmp-inequalities}.

\bibitem{KSTBSA}
J.~Kaniewski, I.~\v{S}upi\'c, J.~Tura, F.~Baccari, A.~Salavrakos and R.~Augusiak, \emph{Maximal nonlocality from
maximal entanglement and mutually unbiased bases, and self-testing of two-qutrit quantum systems}, Quantum
\textbf{3} (2019), 198.

\bibitem{KGZMZ}
D.~Kaszlikowski, P.~Gnaci\'nski, M.~\.Zukowski, W.~Miklaszewski and A.~Zeilinger, \emph{Violations of local realism
by two entangled $N$-dimensional systems are stronger than for two qubits}, Phys. Rev. Lett. \textbf{85} (2000),
4418--4421.

\bibitem{Laeng}
E.~Laeng, \emph{On the $L^p$ norms of the Hilbert transform of a characteristic function}, J. Funct. Anal.
\textbf{262} (2012), no.~10, 4534--4539.

\bibitem{LVN}
B.~Lang, T.~V\'ertesi and M.~Navascu\'es, \emph{Closed sets of correlations: answers from the zoo}, J. Phys. A
\textbf{47} (2014), 424029.

\bibitem{LL}
E.~H. Lieb and M.~Loss, \emph{Analysis}, 2nd ed., Graduate Studies in Mathematics \textbf{14}, American
Mathematical Society, Providence, RI, 2001.

\bibitem{Loomis}
L.~H. Loomis, \emph{A note on the Hilbert transform}, Bull. Amer. Math. Soc. \textbf{52} (1946), 1082--1086.

\bibitem{repo}
A.~Mishra and A.~Senthilkumar, \emph{CGLMP on maximally entangled states: exact optimality, rigidity and
Tsirelson bounds}, preprint and data repository (2026), \url{https://github.com/anshM123/CGLMP}.

\bibitem{MV}
H.~L. Montgomery and R.~C. Vaughan, \emph{Hilbert's inequality}, J. London Math. Soc. (2) \textbf{8} (1974), 73--82.

\bibitem{SATWAP}
A.~Salavrakos, R.~Augusiak, J.~Tura, P.~Wittek, A.~Ac\'{\i}n and S.~Pironio, \emph{Bell inequalities tailored to
maximally entangled states}, Phys. Rev. Lett. \textbf{119} (2017), 040402.

\bibitem{SSKA}
S.~Sarkar, D.~Saha, J.~Kaniewski and R.~Augusiak, \emph{Self-testing quantum systems of arbitrary local dimension
with minimal number of measurements}, npj Quantum Inf. \textbf{7} (2021), 151.

\bibitem{SW}
E.~M. Stein and G.~Weiss, \emph{An extension of a theorem of Marcinkiewicz and some of its applications}, J. Math.
Mech. \textbf{8} (1959), no.~2, 263--284.

\bibitem{TMMD}
X.~K. Tan, J.~Menezes, S.~D. Murthy and A.~C. Dada, \emph{Direct adaptive certification of high-dimensional
entanglement with Bell tests}, arXiv:2608.28235 (2026).

\bibitem{ZG}
S.~Zohren and R.~D. Gill, \emph{Maximal violation of the Collins--Gisin--Linden--Massar--Popescu inequality for
infinite dimensional states}, Phys. Rev. Lett. \textbf{100} (2008), 120406.

\end{thebibliography}

\end{document}
